% Appendix: reasoning points behind the three 7/8 challenge readings, moved from Section 3 (2026-10-27).
% Every factual sentence ends with a comment naming its ledger rows.
\section{Reasoning points for the three 7/8 challenges}\label{app:reasoning}

Tables~\ref{tab:points-QuasiRiemannHypothesis}--\ref{tab:points-HeckeSevenEighths} list the reasoning points behind the QuasiRiemannHypothesis, DirichletSevenEighths and HeckeSevenEighths readings, one ledger row per point; each is INFERRED and marked Read by J. Reilly, 2026-10-27: agree. % v10:L214 v10:L234

\begin{table}[!htbp]
\caption{Reasoning points behind the QuasiRiemannHypothesis reading. INFERRED; Read by J. Reilly, 2026-10-27: agree.}\label{tab:points-QuasiRiemannHypothesis} % v10:L214 v10:L215 v10:L216 v10:L217 v10:L218
\scriptsize\setlength{\tabcolsep}{3pt}\renewcommand{\arraystretch}{0.95}
\begin{tabularx}{\textwidth}{@{}>{\raggedright\arraybackslash}Xl@{}}
\toprule
point & finding \\ % no-fact: table header
\midrule
quantifier order: 7/8 is a fixed literal & MATCH \\ % v10:L214
hypotheses: only Re s > 7/8; no pole hypothesis & MATCH \\ % v10:L215
conclusion: != 0 on the open half-plane Re s > 7/8; the paper excludes the boundary & MATCH \\ % v10:L216
Mathlib meaning: riemannZeta is hurwitzZetaEven 0, the continuation, with a junk value at 1 & MATCH \\ % v10:L217
edge case s = 1: Mathlib's riemannZeta 1 is a nonzero junk value, so the statement asserts a true fact about it and nothing about the pole, which the paper allows & MATCH \\ % v10:L218
\bottomrule
\end{tabularx}
\end{table}

\begin{table}[!htbp]
\caption{Reasoning points behind the DirichletSevenEighths reading. INFERRED; Read by J. Reilly, 2026-10-27: agree.}\label{tab:points-DirichletSevenEighths} % v10:L219 v10:L220 v10:L221 v10:L222 v10:L223 v10:L224
\scriptsize\setlength{\tabcolsep}{3pt}\renewcommand{\arraystretch}{0.95}
\begin{tabularx}{\textwidth}{@{}>{\raggedright\arraybackslash}Xl@{}}
\toprule
point & finding \\ % no-fact: table header
\midrule
quantifier order: fixed 7/8, uniform over every q >= 1, chi and s & MATCH \\ % v10:L219
hypotheses: any modulus q >= 1 and any character, imprimitive ones and q = 1 included, as in the paper's definition & MATCH \\ % v10:L220
hypotheses: hpole excludes s = 1 only for the principal character mod q & MATCH \\ % v10:L221
conclusion: != 0, strict half-plane & MATCH \\ % v10:L222
Mathlib meaning: LFunction is the continuation; for q = 1 it is riemannZeta; for chi != 1 it is entire, so its value at s = 1 is genuine & MATCH \\ % v10:L223
edge cases: imprimitive nonprincipal characters gain zeros only on Re s = 0; the principal pole at 1 is excluded by hpole & MATCH \\ % v10:L224
\bottomrule
\end{tabularx}
\end{table}

\begin{table}[!htbp]
\caption{Reasoning points behind the HeckeSevenEighths reading. INFERRED; Read by J. Reilly, 2026-10-27: agree.}\label{tab:points-HeckeSevenEighths} % v10:L225 v10:L226 v10:L227 v10:L228 v10:L229 v10:L230 v10:L231 v10:L232 v10:L233 v10:L234
\scriptsize\setlength{\tabcolsep}{3pt}\renewcommand{\arraystretch}{0.95}
\begin{tabularx}{\textwidth}{@{}>{\raggedright\arraybackslash}Xl@{}}
\toprule
point & finding \\ % no-fact: table header
\midrule
quantifier order: fixed 7/8, every Character & MATCH \\ % v10:L225
hypotheses: Character (nonzero modulus, MulChar on O/m trivial on units, zero on non-units) is exactly a ray class character mod m extended by zero, since O is a PID with 6 units & MATCH \\ % v10:L226
hypotheses: imprimitive characters are included on both sides & MATCH \\ % v10:L227
hypotheses: the auxiliary period does not matter (HeckeHolo.period\_indep) & MATCH \\ % v10:L228
hypotheses: hpole excludes s = 1 for the principal character only & MATCH \\ % v10:L229
conclusion: != 0, strict half-plane & MATCH \\ % v10:L230
definitions: the lattice series for Re s > 1, holomorphy off \{0, 1\} and the norm formula fix the function in 7/8 < Re s <= 1 by the identity theorem & MATCH \\ % v10:L231
identification: each nonzero ideal has exactly 6 generators; chi and the norm are unit-invariant; so (1/6) times the sum over alpha equals the sum over ideals (absolute convergence) & MATCH \\ % v10:L232
the Dedekind factorization is not needed for this match (checked numerically only) & consistent \\ % v10:L233
edge case s = 0 (outside the half-plane): a Gamma(0) = 0 junk value, no effect on the statement & MATCH \\ % v10:L234
\bottomrule
\end{tabularx}
\end{table}
